\documentclass[11pt]{article}
\usepackage[T1]{fontenc}
\usepackage{lmodern}
\usepackage{amsmath,amssymb,amsthm,mathtools}
\usepackage[margin=1.1in]{geometry}
\usepackage[hidelinks]{hyperref}
\hypersetup{
  pdftitle={Cubic lower bounds for conjugacy separation in nilpotent groups (unverified draft)},
  pdfauthor={Caleb Barnett},
  pdfsubject={Unverified research draft; independent verification is pending.}}
\pagestyle{myheadings}
\markright{\normalfont\small Unverified draft --- 6 October 2026}
\numberwithin{equation}{section}
\newtheorem{theorem}{Theorem}[section]
\newtheorem{proposition}[theorem]{Proposition}
\newtheorem{lemma}[theorem]{Lemma}
\newtheorem{corollary}[theorem]{Corollary}
\theoremstyle{remark}
\newtheorem{remark}[theorem]{Remark}
\newcommand{\Z}{\mathbb Z}
\newcommand{\Q}{\mathbb Q}
\newcommand{\Conj}{\operatorname{Conj}}
\newcommand{\GL}{\operatorname{GL}}
\newcommand{\ad}{\operatorname{ad}}
\newcommand{\ord}{\operatorname{ord}}
\newcommand{\Gp}{\mathcal G}
\title{Cubic lower bounds for conjugacy separation\\in nilpotent groups}
\author{Caleb Barnett}
\date{6 October 2026}

\begin{document}
\maketitle

\begin{center}
\textbf{UNVERIFIED DRAFT}\\[6pt]
\begin{minipage}{0.92\textwidth}
\small
This manuscript presents a proposed proof and is circulated for comment.
The arguments and claimed results have not yet been independently
verified. It was prepared with AI assistance; internal AI-assisted
checks do not establish mathematical correctness. The conclusions
should be treated as provisional. Selected algebraic lemmas have been
formalized in Lean; the main theorem has not. The scope and remaining
formal proof obligations are listed in the
\href{https://github.com/barnettcaleb1/nilpotent-conjugacy}{companion repository}.
Comments and corrections are welcome.
\end{minipage}
\end{center}

\begin{abstract}
This draft proposes a construction of a torsion-free finitely generated
nilpotent group of Hirsch length $8d$ for every integer $d\ge2$, with
a claimed lower bound $n^{3d^2(d+1)}$ for its conjugacy separability
growth. If verified, this would give a negative answer to Question~1
of Der\'e and Vandeputte on the existence of a universal quadratic
bound, in Hirsch length, for the degree of this growth function.
Together with their upper bound, it would show that cubic order in
Hirsch length is necessary and sufficient for a universal degree
bound, up to absolute constants. The argument uses integral
semidirect products, an elementary word compression estimate, and
characters of finite abelian normal subgroups. It addresses arbitrary
finite quotients separating the selected pairs and gives explicit
finite separating quotients. Independent verification of the claims
is pending.
\end{abstract}

\section{Introduction}

Let $G$ be a finitely generated group, and let $S$ be a finite generating
set. We write $|g|_S$ for word length with respect to $S\cup S^{-1}$.
For nonconjugate elements $a,b\in G$, their conjugacy separation depth is
\[
 D_G(a,b)=\min\bigl\{|Q|:\ \varphi:G\twoheadrightarrow Q,\quad
 Q\text{ finite},\quad \varphi(a)\not\sim_Q\varphi(b)\bigr\},
\]
where $\sim_Q$ denotes conjugacy in $Q$. We allow the value $+\infty$
if no such quotient exists. The conjugacy separability growth function is
\[
 \Conj_{G,S}(n)=\max\bigl\{D_G(a,b):\ |a|_S,|b|_S\le n,
 \quad a\not\sim_G b\bigr\},
\]
with value $1$ when the set is empty. Finitely generated nilpotent groups
are conjugacy separable, so this function is finite for the groups under
consideration; see \cite[Introduction]{DV}.

For nondecreasing functions on the positive integers, write $f\preceq g$
if there is an integer $C\ge1$ such that $f(n)\le Cg(Cn)$ for every
$n\ge1$, and write $f\simeq g$ if both comparisons hold. We use
$f\succeq g$ to mean $g\preceq f$. Changing the finite generating set
does not change the resulting growth class, since the corresponding word
lengths are bounded by constant multiples of one another. We therefore
omit $S$ in statements made up to this comparison. All multiplicative
constants may depend on the fixed group and generating set.

Der\'e and Vandeputte prove an upper bound for this function in terms of
the nilpotency class and Hirsch length \cite[Theorem~A]{DV}. They ask
whether there is an absolute constant $C>0$ such that
\begin{equation}\label{eq:question}
 \Conj_G(n)\preceq n^{C h(G)^2}
\end{equation}
for every torsion-free finitely generated nilpotent group $G$
\cite[Question~1, Section~3.2]{DV}. Here $h(G)$ denotes the Hirsch length.
The proposed main result is the following.

\begin{theorem}\label{thm:main}
For every integer $d\ge2$, there is a torsion-free finitely generated
nilpotent group $\Gp_d$ of class $4d$ and Hirsch length $8d$ such that
\[
 \Conj_{\Gp_d}(n)\succeq n^{3d^2(d+1)}.
\]
More precisely, for every finite generating set $S$ of $\Gp_d$, there
are constants $c_{d,S}>0$ and $n_{d,S}\ge1$ such that
\[
 \Conj_{\Gp_d,S}(n)\ge c_{d,S} n^{3d^2(d+1)}
 \qquad\text{for every integer }n\ge n_{d,S}.
\]
Consequently, no absolute constant $C$ satisfies \eqref{eq:question}.
\end{theorem}

The ratio of the exponent in Theorem~\ref{thm:main} to the square of
the Hirsch length is
\[
 \frac{3d^2(d+1)}{(8d)^2}=\frac{3(d+1)}{64},
\]
which tends to infinity. The lower bound is obtained by fixing a single
group $\Gp_d$ and a single prime $p>4d$, and then varying an integer
parameter $s\ge1$ in a sequence of pairs of elements.

If verified, the claimed lower bound would also determine the order of
the optimal universal exponent bound as a function of Hirsch length.

\begin{corollary}\label{cor:envelope}
Every torsion-free finitely generated nilpotent group $G$ satisfies
\[
 \Conj_G(n)\preceq n^{h(G)^3}.
\]
There is an absolute constant $c_0>0$ and a sequence of such groups
$G$ with unbounded Hirsch length for which
\[
 \Conj_G(n)\succeq n^{c_0 h(G)^3}.
\]
In particular, no function $E:\mathbb N_0\to[0,\infty)$ with
$E(h)=o(h^3)$ can satisfy $\Conj_G(n)\preceq n^{E(h(G))}$ for all
torsion-free finitely generated nilpotent groups $G$.
\end{corollary}

\begin{proof}
Suppose first that $G$ is nonabelian. Put $h=h(G)$, let $k$ be its
nilpotency class, and put $z=h(Z(G))$ and $r_0=h([G,G])$.
By \cite[Theorem~A]{DV},
\[
 \Conj_G(n)\preceq n^{(k-1)(h-z+1)r_0}.
\]
The inequalities $k-1\le h$, $z\ge1$, and $r_0\le h$ bound this exponent
by $h^3$. If $G$ is nontrivial abelian, its growth is equivalent to
$\log n$ by \cite[Proposition~2.1.9]{DV}; the trivial group has constant
growth. This proves the upper bound. For the lower bound, take
$G=\Gp_d$ and $c_0=3/512$, since
\[
 \frac{3d^2(d+1)}{(8d)^3}=\frac{3(d+1)}{512d}\ge\frac{3}{512}.
\]
For any proposed function $E=o(h^3)$, choose one sufficiently large
$d$ so that $E(8d)<c_0(8d)^3$. The two polynomial bounds in this fixed
group are then incompatible as $n\to\infty$.
\end{proof}

The groups in the proof have the form $G(m)=\Z^m\rtimes H$, where
$H$ is itself an abelian-by-cyclic nilpotent group. Two actions on the
normal subgroup $\Z^m$ play different roles. One produces a bidiagonal
conjugacy equation; the other forces sufficiently many matrix units to
act on the abelian subgroup in every finite quotient. A character
detecting the chosen conjugacy obstruction then gives many independent
characters of large order. Section~\ref{sec:compression} supplies the
word estimates, and Sections~\ref{sec:group}--\ref{sec:finish} give the
construction and quotient argument.

Throughout, group commutators are $[x,y]=xyx^{-1}y^{-1}$, vectors are
columns, and exponentials and logarithms of nilpotent or unipotent
matrices are finite polynomial expressions.

\section{An elementary word compression estimate}\label{sec:compression}

We first prove the word estimate needed in the construction. All word
length comparisons in this section are upper bounds.

\begin{lemma}\label{lem:blockcompression}
Let $l\ge1$ be an integer and let $F_l=\Z^l\rtimes_A\Z$, where $A=I+N$ and
\[
 Ne_i=e_{i+1}\quad(1\le i<l),\qquad Ne_l=0.
\]
Use the generating set $e_1,\ldots,e_l,t$, where $t$ generates the cyclic
factor. For every $1\le j\le l$ and every $a\in\Z$,
\[
 |ae_j|_{F_l}\le C_l(1+|a|^{1/j}).
\]
\end{lemma}

\begin{proof}
We use descending induction on $j$. It suffices to consider $a>0$:
inversion handles negative $a$, and $a=0$ is immediate. Put
$B=\lfloor a^{1/j}\rfloor+1\ge2$. Since $a<B^j$, write
\[
 a=\sum_{u=0}^{j-1}a_u B^u,\qquad 0\le a_u<B.
\]
The commutator identity
\[
 [t^v,x]=(A^v-I)x\qquad(x\in\Z^l)
\]
shows that iterated commutators with $u$ copies of $t^B$ and
$j-1-u$ copies of $t$, applied to $a_ue_1$, represent
\[
 x_u=a_u(A^B-I)^uN^{j-1-u}e_1.
\]
The operators commute, so their order in the iteration is immaterial.
Each such word has length $O_l(B)$: its depth is at most $l-1$, and
every input word has length at most $B$. When $j=1$, this simply means
the word representing $a_0e_1$.

The first nonzero possible coordinate of $x_u$ is its $e_j$ coordinate,
which is $a_u B^u$. To control the higher coordinates, expand
\[
 A^B-I=\sum_{v=1}^{l-1}\binom Bv N^v.
\]
For a contribution to the $e_i$ coordinate with $i>j$, the indices from
the $u$ factors sum to $i-j+u$. Its coefficient is therefore bounded by
$C_lB^{1+i-j+u}$, including the digit $a_u$. As $u\le j-1$, this is
at most $C_lB^i$. Consequently the sum of the $j$ vectors is
\[
 \sum_{u=0}^{j-1}x_u=ae_j+\sum_{i>j}b_i e_i,
 \qquad |b_i|\le C_lB^i,
\]
and has a representative of length $O_l(B)$.

For $j=l$ there are no error coordinates, which starts the induction.
For smaller $j$, the induction hypothesis represents every $b_i e_i$
with $i>j$ by a word of length $O_l(B)$. Cancel these finitely many
vectors. They all lie in the abelian subgroup $\Z^l$, so the
cancellations introduce no further errors. This gives a word of length
$O_l(B)$ for $ae_j$, as required.
\end{proof}

\begin{lemma}\label{lem:compression}
Let $l\ge1$ be an integer and let $K=\Z^l\rtimes_A\Z$, where $A\in\GL_l(\Z)$ is unipotent and
$N=A-I$ has a single Jordan block over $\Q$. If
$v\in N^{j-1}\Q^l\cap\Z^l$, where $1\le j\le l$, then
\[
 |av|_K\le C_{A,v,S}(1+|a|^{1/j})\qquad(a\in\Z)
\]
for every fixed finite generating set $S$ of $K$.
\end{lemma}

\begin{proof}
Choose an integral cyclic vector $f_1$ for $N$ and set
$f_i=N^{i-1}f_1$. These vectors form a rational basis and span a
finite-index sublattice $F\subset\Z^l$ preserved by $A$ and $A^{-1}$.
The action on this basis is exactly the action in
Lemma~\ref{lem:blockcompression}. The subgroup $F\rtimes_A\Z$ thus
has words for $af_i$ of length $O_l(1+|a|^{1/i})$ in its indicated
generators. Each of these generators has bounded $S$-length in $K$,
so the same upper estimates hold in $K$ with a different constant.

The vectors $f_j,\ldots,f_l$ span $N^{j-1}\Q^l$. Hence there is an
integer $D\ge1$ such that
$Dv=\sum_{i=j}^l c_i f_i$ with $c_i\in\Z$. Write
$a=Da_0+b$, where $0\le b<D$. Then
\[
 av=\sum_{i=j}^l a_0c_i f_i+bv.
\]
Apply the estimates just proved to the terms in the sum, and absorb the
finitely many possibilities for $bv$ into the constant. Since $i\ge j$,
each term is bounded by a constant times $1+|a|^{1/j}$.
\end{proof}

\section{The integral semidirect product}\label{sec:group}

Fix an integer $m\ge4$. Let $V=\Z^m$, with basis $e_1,\ldots,e_m$,
and define endomorphisms $J,Y$ of $V$ by
\[
 Je_i=e_{i+1},\qquad Ye_i=i e_{i+1}\quad(1\le i<m),
 \qquad Je_m=Ye_m=0.
\]
Their matrix commutators satisfy
\begin{equation}\label{eq:commutator}
 [Y,J^j]=YJ^j-J^jY=jJ^{j+1}.
\end{equation}
Put $M=m!$ and consider the additive group of commuting matrices
\[
 W=M\Z J\oplus M\Z J^2\oplus\cdots\oplus M\Z J^{m-1}.
\]
Conjugation by $\exp Y$ induces an automorphism $T$ of $W$. Indeed,
iterating \eqref{eq:commutator} gives
\begin{equation}\label{eq:T}
 T(J^j)=\exp(\ad Y)J^j=J^j(1-J)^{-j},\qquad
 T^{-1}(J^j)=J^j(1+J)^{-j}.
\end{equation}
Both expressions are finite integral series, since $J^m=0$.

Define
\[
 H=W\rtimes_T\Z,\qquad
 (P,a)(Q,b)=(P+T^aQ,a+b),
\]
and set
\begin{equation}\label{eq:rho}
 \rho(P,a)=\exp(P)\exp(aY).
\end{equation}
The conjugation identity defining $T$ shows that $\rho$ is a
representation. It takes values in $\GL_m(\Z)$. For $P\in W$, both
$\exp(P)$ and $\exp(-P)$ are integral: $P$ is divisible by $M=m!$,
which clears all factorial denominators in these finite series. Also,
\begin{equation}\label{eq:expY}
 \exp(aY)e_i=\sum_{v=0}^{m-i}a^v\binom{i+v-1}{v}e_{i+v}
\end{equation}
is integral for every $a\in\Z$, including negative $a$.

Our group is
\begin{equation}\label{eq:Gm}
 G(m)=V\rtimes_\rho H,\qquad
 (v,h)(v',h')=(v+\rho(h)v',hh').
\end{equation}
We identify $V$ and $H$ with their natural subgroups of $G(m)$.

\begin{proposition}\label{prop:group}
The group $G(m)$ is torsion-free and finitely generated, has Hirsch
length $2m$, and is nilpotent of class $m$. The vector $e_m$ is central.
\end{proposition}

\begin{proof}
The successive extensions by the free abelian groups $V$, $W$, and
$\Z$ prove finite generation and torsion-freeness. Additivity of Hirsch
length in these extensions gives
\[
 h(G(m))=m+(m-1)+1=2m.
\]
To see nilpotency, first note that $\rho$ is faithful. If
$\exp(P)\exp(aY)=I$, then $\exp(P)=\exp(-aY)$. Taking finite
polynomial logarithms gives $P=-aY$. The first lower diagonal of $P$
has a constant entry, whereas that of $Y$ has entries $1,\ldots,m-1$.
As $m\ge4$, this forces $a=0$ and then $P=0$.

It follows that
\begin{equation}\label{eq:affine}
 (v,h)\longmapsto
 \begin{pmatrix}1&0\\v&\rho(h)\end{pmatrix}
\end{equation}
is an embedding into the lower unitriangular integer matrices of size
$m+1$. Thus $G(m)$ has class at most $m$. Let $z=(MJ,0)\in H$.
The iterated commutator with $m-1$ copies of $z$ and one copy of $e_1$
is the vector
\[
 (\exp(MJ)-I)^{m-1}e_1=M^{m-1}e_m\ne0.
\]
It lies in the $m$th term of the ordinary lower central series, proving
that the class is exactly $m$. Finally, every $\rho(h)$ fixes $e_m$,
so this vector is central in $G(m)$.
\end{proof}

Fix the generating set consisting of $e_1,\ldots,e_m$, the elements
$(MJ^j,0)\in H$ for $1\le j<m$, and the stable letter
$u=(0,1)\in H$. Until the final step, word lengths refer to this set
or to its indicated subsets in the relevant subgroups.

\begin{lemma}\label{lem:groupcompression}
For integers $w,t$ with $1\le w\le m-2$ and $t\ge1$,
\begin{equation}\label{eq:actinglength}
 |(M(tJ^w+J^{w+1}),0)|_H\le C_m(1+t^{1/w}).
\end{equation}
For every $a\in\Z$,
\begin{equation}\label{eq:centrallength}
 |aM^{m-1}e_m|_{G(m)}\le C_m(1+|a|^{1/m}).
\end{equation}
\end{lemma}

\begin{proof}
On $W$, the operator $T-I$ has one rational Jordan block. Indeed,
\eqref{eq:T} gives
\[
 (T-I)(MJ^j)=jMJ^{j+1}+\text{higher powers of }J.
\]
The chain starting at $MJ$ therefore has successive leading terms
$M(i-1)!J^i$, for $1\le i<m$. In particular, $MJ^w$ lies in the
rational image of $(T-I)^{w-1}$. Apply Lemma~\ref{lem:compression}
to $H=W\rtimes_T\Z$ and add the fixed vector $MJ^{w+1}$ to obtain
\eqref{eq:actinglength}.

For \eqref{eq:centrallength}, use the subgroup generated by $V$ and
$z=(MJ,0)\in H$. Its action on $V$ is $\exp(MJ)$, and the difference
$\exp(MJ)-I$ has one rational Jordan block. Its chain beginning at
$e_1$ ends in $M^{m-1}e_m$. Lemma~\ref{lem:compression}, with $j=m$,
gives the desired words in this subgroup, whose generators have bounded
length in $G(m)$.
\end{proof}

\section{Short nonconjugate pairs}\label{sec:pairs}

From now on fix integers and a prime satisfying
\begin{equation}\label{eq:parameters}
 m\ge4,\qquad
 L=\left\lceil\frac{m-1}{2}\right\rceil\le w\le m-2,
 \qquad p>m,
\end{equation}
and put $r=m-w$. The group $G(m)$ and the prime $p$ remain fixed.
For each integer $s\ge1$, define
\begin{equation}\label{eq:pairs}
 \begin{gathered}
 t=p^s,\qquad P_s=MJ^w(tI+J),\\
 a_s=(0,(P_s,0)),\qquad c_s=M^{m-1}p^{s-1}e_m,\qquad
 b_s=a_sc_s.
 \end{gathered}
\end{equation}
Here $c_s$ is viewed as a central element of $G(m)$.

\begin{proposition}\label{prop:pairs}
The elements $a_s$ and $b_s$ are nonconjugate in $G(m)$ and satisfy
\begin{equation}\label{eq:pairlength}
 |a_s|,|b_s|\le C_{m,p}p^{s/w}\qquad(s\ge1).
\end{equation}
\end{proposition}

\begin{proof}
Lemma~\ref{lem:groupcompression} bounds $|a_s|$ by
$C_m(1+p^{s/w})$ and $|c_s|$ by $C_m(1+p^{(s-1)/m})$.
Since $w<m$ and $b_s=a_sc_s$, these bounds imply
\eqref{eq:pairlength} after enlarging the constant.

For nonconjugacy, write
\begin{equation}\label{eq:R}
 R_s=\exp(P_s)-I=P_sB_s,\qquad
 B_s=\sum_{v=0}^{m-1}\frac{P_s^v}{(v+1)!}.
\end{equation}
Every term of $B_s$ is integral because $P_s$ is divisible by $m!$.
Moreover, $B_s-I$ is nilpotent, so its finite geometric inverse is
integral. Thus $B_s\in\GL_m(\Z)$ and
\begin{equation}\label{eq:integralimage}
 R_sV=P_sV.
\end{equation}

Put $h_s=(P_s,0)\in H$. If $(v,k)\in G(m)$ conjugates $a_s$ to
$b_s$, equality in the $H$ coordinate requires
$kh_sk^{-1}=h_s$. Under this condition the translation coordinate of
$(v,k)a_s(v,k)^{-1}$ is $(I-\rho(h_s))v=-R_sv$. Conversely, any
vector in $R_sV$ is obtained in this way by a conjugator in $V$.
Therefore
\begin{equation}\label{eq:conjequivalence}
 a_s\sim_{G(m)}b_s\quad\Longleftrightarrow\quad
 c_s\in R_sV=P_sV.
\end{equation}

Let $f=J^w(tI+J)$. Its nonzero square block, from the input
coordinates $e_1,\ldots,e_r$ to the output coordinates
$e_{w+1},\ldots,e_m$, is
\[
 \begin{pmatrix}
 t&0&\cdots&0\\
 1&t&\ddots&\vdots\\
 0&\ddots&\ddots&0\\
 0&\cdots&1&t
 \end{pmatrix}.
\]
Solving $fx=k e_m$ over $\Z$ successively gives
$x_1=\cdots=x_{r-1}=0$ and $t x_r=k$. Hence $k e_m\in fV$ if
and only if $t$ divides $k$. Since $P_s=Mf$, membership of $c_s$ in
$P_sV$ would require
\[
 p^s\mid M^{m-2}p^{s-1}.
\]
This is impossible because $p>m$ and $M=m!$. Equation
\eqref{eq:conjequivalence} proves the assertion.
\end{proof}

\section{Lower bounds for arbitrary finite quotients}\label{sec:quotients}

Let $\varphi:G(m)\twoheadrightarrow Q$ be any finite quotient in which
$\varphi(a_s)$ and $\varphi(b_s)$ are nonconjugate. We shall bound
$|Q|$ by examining the finite abelian normal subgroup
$A=\varphi(V)$.

\subsection{The obstruction in the \texorpdfstring{$p$}{p}-primary component}

Conjugation by $\varphi(a_s)$ acts on $A$ as $I+R_s$, so its
difference from the identity acts as $R_s$. If $\varphi(c_s)$ belonged
to $R_sA$, an element of $A$ would conjugate $\varphi(a_s)$ to
$\varphi(b_s)$. Consequently,
\begin{equation}\label{eq:Aobstruction}
 \varphi(c_s)\notin R_sA.
\end{equation}
On the other hand, there is an exact identity in $V$:
\begin{equation}\label{eq:pc}
 R_s(M^{m-2}e_{m-w})=M^{m-1}p^s e_m=p c_s.
\end{equation}
To check it, the first application of $P_s$ gives the displayed
multiple of $e_m$, and every further application is zero.

Let $A_p$ be the $p$-primary component of $A$, and let
$\pi:V\twoheadrightarrow A_p$ be the composition of $\varphi|_V$
with the primary projection. Every endomorphism of $A$ preserves its
primary components. On a component at a prime other than $p$,
multiplication by $p$ is invertible. Equation~\eqref{eq:pc} thus puts
the corresponding component of $\varphi(c_s)$ in the image of $R_s$.
It follows from \eqref{eq:Aobstruction} that
\begin{equation}\label{eq:pobstruction}
 \pi(c_s)\notin R_sA_p.
\end{equation}

For the remaining module calculations, use the ring
\[
 \mathcal R=\Z_{(p)}
 =\{a/b\in\Q:\ a,b\in\Z,\ p\nmid b\}.
\]
Any finite abelian $p$-group is naturally an $\mathcal R$-module,
because multiplication by $b$ is invertible whenever $p\nmid b$.
The map $\pi$ extends uniquely to a surjective $\mathcal R$-linear map
$\mathcal R^m\to A_p$. Write $\Lambda$ for its kernel.

\begin{lemma}\label{lem:descent}
The submodule $\Lambda$ is invariant under $J$ and $Y$. These
operators therefore induce endomorphisms of $A_p$, and
\begin{equation}\label{eq:fobstruction}
 \pi(c_s)\notin fA_p,\qquad f=J^w(tI+J).
\end{equation}
\end{lemma}

\begin{proof}
The matrices $U=\exp(MJ)$ and $E=\exp(Y)$ are the actions of
$z=(MJ,0)$ and $u=(0,1)$ in $G(m)$. The kernel $\Lambda$ is
therefore invariant under both matrices, and hence under $U-I$ and
$E-I$. Since $(U-I)^m=(E-I)^m=0$, the finite identities
\[
 \log U=\sum_{v=1}^{m-1}\frac{(-1)^{v+1}}v(U-I)^v=MJ,
 \qquad
 \log E=\sum_{v=1}^{m-1}\frac{(-1)^{v+1}}v(E-I)^v=Y
\]
show that $\Lambda$ is invariant under $MJ$ and $Y$: every
denominator is a unit in $\mathcal R$, because $p>m$.
As $M$ is also a unit in this ring, $\Lambda$ is invariant under $J$.

The matrices $B_s$ and $B_s^{-1}$ in \eqref{eq:R} are polynomials
in $J$ with coefficients in $\mathcal R$. They preserve $\Lambda$
and induce inverse automorphisms of $A_p$. Multiplication by $M$ is
an automorphism as well. Thus $R_sA_p=P_sA_p=fA_p$, and
\eqref{eq:pobstruction} gives \eqref{eq:fobstruction}.
\end{proof}

\subsection{Matrix units forced by the two actions}

For $1\le i,j\le m$, let $E_{j,i}$ denote the matrix sending $e_i$
to $e_j$ and all other basis vectors to zero.

\begin{lemma}\label{lem:matrixunits}
If $j-i\ge L=\lceil(m-1)/2\rceil$, then $E_{j,i}$ belongs to the
associative $\mathcal R$-algebra generated by $J$ and $Y$. In
particular, it preserves $\Lambda$ and induces an endomorphism of
$A_p$.
\end{lemma}

\begin{proof}
Fix a gap $\ell\ge L$. For $0\le v\le\ell$, the operator
$J^{\ell-v}Y^v$ sends $e_i$, for $1\le i\le m-\ell$, to
\[
 i(i+1)\cdots(i+v-1)e_{i+\ell}.
\]
For $v=0$ the coefficient is $1$. The coefficient polynomials are
monic of degrees $0,\ldots,\ell$, so form an $\mathcal R$-basis
for the polynomials of degree at most $\ell$.

There are $m-\ell$ active input indices, and
$m-\ell-1\le\ell$. Given one of these indices $i_0$, the Lagrange
polynomial
\[
 \prod_{\substack{1\le i\le m-\ell\\i\ne i_0}}
 \frac{X-i}{i_0-i}
\]
has degree at most $\ell$ and takes the value $1$ at $i_0$ and $0$
at every other active index. Its denominators are units in
$\mathcal R$, since all their nonzero factors have absolute value
less than $m<p$. Expressing this polynomial in the preceding basis
gives an $\mathcal R$-linear combination of the operators
$J^{\ell-v}Y^v$ that is exactly $E_{i_0+\ell,i_0}$. Every inactive
input index is killed by each of these operators. The final assertion
follows from Lemma~\ref{lem:descent}.
\end{proof}

\subsection{Independent characters of large order}

\begin{proposition}\label{prop:quotientbound}
Under \eqref{eq:parameters}, every finite quotient $Q$ separating
$a_s$ and $b_s$ satisfies
\begin{equation}\label{eq:quotientbound}
 |Q|\ge p^{s(m-w)(w+1-L)}.
\end{equation}
\end{proposition}

\begin{proof}
By \eqref{eq:fobstruction}, $\pi(c_s)$ is nonzero in the finite
abelian group $A_p/fA_p$. Characters into $\Q/\Z$ separate points
of a finite abelian group: this follows by decomposing the group into
cyclic summands and using the standard characters on those summands.
We can therefore choose
\[
 \lambda:A_p\longrightarrow\Q/\Z
\]
which vanishes on $fA_p$ and satisfies $\lambda(\pi(c_s))\ne0$.
Write $\bar e_i=\pi(e_i)$, $\lambda_i=\lambda(\bar e_i)$, and
$\alpha=\lambda_{w+1}$.

The successive nonzero columns of $f=J^w(tI+J)$ give
\[
 t\lambda_{w+i}+\lambda_{w+i+1}=0\quad(1\le i<r),
 \qquad t\lambda_m=0.
\]
Thus
\begin{equation}\label{eq:recurrence}
 \lambda_{w+i}=(-t)^{i-1}\alpha\quad(1\le i\le r).
\end{equation}
Since $c_s=M^{m-1}p^{s-1}e_m$ and $M$ is prime to $p$, the
nonzero value of $\lambda(\pi(c_s))$, together with
$p^s\lambda_m=0$, shows that $\lambda_m$ has exactly order $p^s$.
Equation~\eqref{eq:recurrence} then implies
\begin{equation}\label{eq:alphaorder}
 \ord(\alpha)=p^{sr}.
\end{equation}
Indeed, $\lambda_m=(-p^s)^{r-1}\alpha$, and every value of a
character on $A_p$ has order a power of $p$.

Put $q=w+1-L$. For $1\le i\le q$, the gap $w+1-i$ is at least
$L$, so Lemma~\ref{lem:matrixunits} makes
\[
 \chi_i=\lambda\circ E_{w+1,i}
\]
a well-defined character of $A_p$. Its values on the generating
elements $\bar e_j$ are
\[
 \chi_i(\bar e_j)=
 \begin{cases}\alpha,&j=i,\\0,&j\ne i.\end{cases}
\]
In particular, each $\chi_i$ has exactly order $p^{sr}$. If
$\sum_{i=1}^q b_i\chi_i=0$, evaluating on $\bar e_i$ gives
$b_i\alpha=0$, hence $p^{sr}\mid b_i$ for every $i$. The characters
$\chi_1,\ldots,\chi_q$ therefore generate a subgroup isomorphic to
$(\Z/p^{sr}\Z)^q$ in the character group of $A_p$. A finite
abelian group and its character group have the same order, so
\[
 |Q|\ge |A_p|\ge p^{srq}.
\]
Since $r=m-w$ and $q=w+1-L$, this is \eqref{eq:quotientbound}.
\end{proof}

\section{Explicit finite separating quotients}\label{sec:detectors}

The preceding lower bound applies to every finite separating quotient.
For completeness, the same matrices give a finite separating quotient
for each selected pair directly.

\begin{proposition}\label{prop:detector}
The image of the affine representation \eqref{eq:affine} modulo
$p^{sr}$ is a finite quotient of $G(m)$ in which $a_s$ and $b_s$
remain nonconjugate.
\end{proposition}

\begin{proof}
On $V/p^{sr}V$, define an additive character $\mu$ by
\[
 \mu(e_i)=0\quad(1\le i\le w),\qquad
 \mu(e_{w+i})=\frac{(-p^s)^{i-1}}{p^{sr}}+\Z
 \quad(1\le i\le r).
\]
These values are annihilated by $p^{sr}$, so the character is
well-defined. The adjacent columns of $f$ give cancellation, and its
last nonzero column gives
\[
 p^s\mu(e_m)=(-1)^{r-1}+\Z=0.
\]
Hence $\mu$ vanishes on $fV$, and therefore on
$R_sV=P_sB_sV=MfB_sV$. In contrast,
\begin{equation}\label{eq:mudetects}
 \mu(c_s)=\frac{(-1)^{r-1}M^{m-1}}p+\Z\ne0.
\end{equation}

Write $E_s=\exp(P_s)$. The matrices of $a_s$ and $b_s$ are
respectively
\[
 \begin{pmatrix}1&0\\0&E_s\end{pmatrix},\qquad
 \begin{pmatrix}1&0\\c_s&E_s\end{pmatrix}.
\]
If their reductions were conjugate in the finite affine image, a
conjugating matrix would have the form
$\left(\begin{smallmatrix}1&0\\v&K\end{smallmatrix}\right)$.
Equality of the lower-right blocks would require
$KE_sK^{-1}=E_s$ modulo $p^{sr}$. Equality of the translation
columns would then require
\[
 c_s=(I-E_s)v=-R_sv\pmod{p^{sr}V}.
\]
Applying $\mu$ contradicts \eqref{eq:mudetects}. Thus this finite
matrix image separates the pair.
\end{proof}

\section{Proof of the main theorem}\label{sec:finish}

Fix $d\ge2$, take $m=4d$ and $w=3d$, and set $\Gp_d=G(4d)$.
These choices satisfy \eqref{eq:parameters}: here
\[
 L=2d,\qquad r=d,\qquad w+1-L=d+1.
\]
Proposition~\ref{prop:group} gives class $4d$ and Hirsch length
$8d$. Choose one prime $p>4d$ and keep it fixed. For every $s\ge1$,
Propositions~\ref{prop:pairs}, \ref{prop:quotientbound}, and
\ref{prop:detector} give nonconjugate elements with
\begin{equation}\label{eq:finalpairs}
 |a_s|,|b_s|\le C p^{s/(3d)},\qquad
 p^{sd(d+1)}\le D_{\Gp_d}(a_s,b_s)<\infty,
\end{equation}
where $C\ge1$ is independent of $s$. For any other fixed finite
generating set $S$, enlarge $C$ to preserve the first inequality.

Let $n\ge C p^{1/(3d)}$, and choose the largest integer $s\ge1$
such that $C p^{s/(3d)}\le n$. Maximality gives
$n<Cp^{(s+1)/(3d)}$. Raising this inequality to the power
$E_d=3d^2(d+1)$ and using \eqref{eq:finalpairs}, we obtain
\[
 \Conj_{\Gp_d,S}(n)
 \ge p^{sd(d+1)}
 > C^{-E_d}p^{-d(d+1)}n^{E_d}.
\]
This proves the asserted bound for all sufficiently large integers
$n$, with constants allowed to depend on $d$ and $S$.

Finally, suppose an absolute $C_0$ satisfied
$\Conj_G(n)\preceq n^{C_0h(G)^2}$ for every group in question.
Choose $d$ so that $3(d+1)/64>C_0$. In this one fixed group
$\Gp_d$, the lower exponent $3d^2(d+1)$ is strictly greater than
$C_0(8d)^2$. The definition of $\preceq$ would bound
$\Conj_{\Gp_d,S}(n)$ by a constant multiple of
$n^{C_0(8d)^2}$, contradicting the lower bound as $n\to\infty$.
This completes the proof of Theorem~\ref{thm:main}.

\begin{remark}
The source of the exponent is visible in the bidiagonal conjugacy
equation. A character annihilating $fA_p$ satisfies
$\lambda_m=(-p^s)^{r-1}\lambda_{w+1}$. Detecting the chosen multiple
$p^{s-1}e_m$ forces $\lambda_m$ to have order $p^s$, and consequently
$\lambda_{w+1}$ to have order $p^{sr}$. The second action $Y$,
together with $J$, supplies the matrix units that reproduce this
value in $w+1-L$ independent characters. Combining their orders with
the word estimate of weight $w$ yields the exponent
$w(m-w)(w+1-L)$. Theorem~\ref{thm:main} uses
$m=4d$ and $w=3d$; it does not assert an exact growth exponent for
$\Gp_d$.
\end{remark}

\begin{thebibliography}{9}
\bibitem{DV}
Jonas Der\'e and Lukas Vandeputte,
\emph{Uniform Upper Bounds for Conjugacy Separability Growth in
Nilpotent Groups},
arXiv:2610.03074v1 [math.GR], 2 October 2026.
\href{https://arxiv.org/abs/2610.03074v1}{\texttt{arxiv.org/abs/2610.03074v1}}.
\end{thebibliography}

\end{document}
